\section{Menu 3: single-photoelectron calibration}

\begin{frame}{Menu 3: goal}
\begin{itemize}
\item The peak detector, after the slow shaper, follows the signal from the opening of the
  window, holds the highest value reached and stops at the hold; the ADC reads that value.
\item The simulation expresses every charge in photoelectrons; the chip stores ADC counts.
  The gain of a channel, in counts per photoelectron, is the unit of the whole chain.
\item Above breakdown, with a periodic trigger, the window opens at a random time: it holds
  the pedestal or a dark pulse of 1, 2, 3~pe. The spectrum is a ladder of equally spaced
  peaks; the spacing is the gain.
\item Measured for every channel of every board: the gain at three biases and three shaping
  times; the breakdown voltage from the gain against the bias; the pedestal noise of menu~1 in
  photoelectrons.
\item Staircases at two biases: one photoelectron in units of the threshold.
\end{itemize}
\end{frame}

\begin{frame}{Menu 3: data and method}
\begin{columns}
\column{0.5\textwidth}
\small
\begin{tabular}{llll}
\toprule
Run & Bias & HG shaping & Boards \\
\midrule
108--110 & 43 V & 25, 50, 87.5 ns & 1, 2, 4 \\
111 & 42 V & 25 ns & 1, 2, 4 \\
112 & 41 V & 25 ns & 1, 2, 4 \\
\midrule
Scan & Bias & Thresholds & Board \\
\midrule
113 & 43 V & 184--385 & 1 \\
114, 115 & 41 V & 150--300, 190--275 & 1 \\
\bottomrule
\end{tabular}

\medskip
Periodic trigger at 1~kHz, HoldDelay 200~ns, preamplifier gains 55, peak-stretch mode,
50\,000 events per board. Bias at the SiPM: HV output minus 4.2~V.
\column{0.5\textwidth}
\begin{itemize}
\item The window opens at a random time: the spectrum is a ladder of 0, 1, 2, \ldots
  photoelectrons.
\item Each channel is fitted with a comb of six Gaussians at $p + n g$, widths
  $\sqrt{\sigma_0^2 + n\sigma_1^2}$: $g$ is the gain in counts per photoelectron.
\item Resolved when the 1-pe peak is narrower than half the gain.
\item Breakdown voltage: intercept of the line through the gain at 41, 42 and 43~V.
\item The six channels without SiPM signal are left out.
\end{itemize}
\end{columns}
\end{frame}

\begin{frame}{The ladder of one channel}
\centering
\includegraphics[width=0.88\textwidth]{spe_spectra.png}
\end{frame}

\begin{frame}{Gain of every channel}
\begin{columns}
\column{0.62\textwidth}
\includegraphics[width=\textwidth]{spe_gains.png}
\column{0.38\textwidth}
\begin{itemize}
\item At 25~ns the ladder is resolved in 183, 184 and 161 of the 186 channels at 41, 42 and
  43~V.
\item Gain, medians of the three boards: 85--92 counts per photoelectron at 41~V, 127--134 at
  42~V, 165--178 at 43~V.
\item Spread across the channels of a board: 2.4--5.6\%.
\item Width of the 1-pe peak: 25--30 counts at 41~V, 49--52 at 43~V.
\end{itemize}
\end{columns}
\end{frame}

\begin{frame}{Breakdown voltage}
\begin{columns}
\column{0.62\textwidth}
\includegraphics[width=\textwidth]{spe_breakdown.png}
\column{0.38\textwidth}
\begin{itemize}
\item Linear in the bias: 41--42 counts per photoelectron per volt.
\item Breakdown 38.8, 38.8, 39.0~V (boards 1, 2, 4), rms 0.25--0.35~V over the channels.
\item Overvoltage 4.1, 3.1, 2.1~V at 43, 42, 41~V.
\item Expected: breakdown 38~V (43~V set as 5~V of overvoltage). Observed: 0.8--1.0~V
  higher.
\end{itemize}
\end{columns}
\end{frame}

\begin{frame}{Gain against shaping}
\begin{columns}
\column{0.58\textwidth}
\includegraphics[width=\textwidth]{spe_shaping.png}
\column{0.42\textwidth}
\begin{itemize}
\item At 43~V the ladder is resolved in 110 channels at 50~ns and 26 at 87.5~ns (1-pe width
  61--67 counts).
\item Gain: 165--178 at 25~ns, 176--182 at 50~ns, 170--173 at 87.5~ns.
\item Menu 2 (stars) and menu 3 agree within 12 counts in the three channels of board 1 where
  the ladder is resolved.
\end{itemize}
\end{columns}
\end{frame}

\begin{frame}{Pedestal noise in photoelectrons}
\centering
\scriptsize
\input{\results/spe_noise_pe.tex}

\medskip
\normalsize
\begin{itemize}
\item At 43~V and 25~ns: 0.068--0.075 photoelectrons per channel, 0.033--0.038 coherent,
  0.059--0.065 incoherent.
\item At 50~ns: 0.067--0.074, coherent 0.041--0.046.
\item Expected (model of the simulation at 4.1~V of overvoltage): 0.137, 0.097, 0.073
  photoelectrons at 25, 50, 87.5~ns. Observed: 0.07 at the three shapings.
\end{itemize}
\end{frame}

\begin{frame}{Staircases at 5 and 3 V, and the entries at zero}
\centering
\scriptsize
\input{\results/spe_staircases.tex}

\medskip
\input{\results/spe_zeros.tex}

\medskip
\normalsize
\begin{itemize}
\item One photoelectron is 22 DAC units at 43~V and 12.5 at 41~V (gain ratio 1.8); one DAC
  unit is 7.8 and 7.1 counts of the high gain.
\item Entries at 0 at 43~V: 0.5--1.1\% at 25~ns, 3.8--6.1\% at 50~ns, 8.9--12\% at 87.5~ns;
  none at 41~V. Open (F12).
\end{itemize}
\end{frame}

\begin{frame}{Menu 3: summary of results}
\begin{itemize}
\item At 25~ns the ladder is resolved in nearly every channel (183, 184, 161 of 186 at 41,
  42, 43~V); at 50~ns in 110; at 87.5~ns in 26, where the dark pulses pile up.
\item The gain is linear in the bias, 41--42 counts per photoelectron per volt; breakdown at
  38.8--39.0~V on the three boards, uniform within 0.3~V across the channels. The bias of
  43~V is 4.1~V above breakdown, not 5.
\item The gain varies by 2.4--5.6\% across the channels of a board and little with the
  shaping.
\item The pedestal noise of the chain at 43~V is 0.07 photoelectrons per channel at the three
  shapings, half of it coherent; the model has 0.073 at 87.5~ns and 0.137 at 25~ns.
\item One DAC unit of the threshold is 7.8 counts of the high gain at 43~V.
\item Open: at 43~V a fraction of the entries, 0.5 to 12\% growing with the shaping, reads 0
  (F12).
\end{itemize}
\end{frame}

\begin{frame}{Menu 3: values}
\begin{enumerate}
\item Gain at 25~ns: 85--92 counts per photoelectron at 41~V, 127--134 at 42~V, 165--178 at
  43~V; spread across the channels of a board 2.4--5.6\%.
\item Breakdown voltage 38.8, 38.8, 39.0~V (boards 1, 2, 4), rms 0.25--0.35~V; 41--42 counts
  per photoelectron per volt; 43~V is 4.0--4.2~V above breakdown.
\item Gain at 43~V: 165--178 at 25~ns, 176--182 at 50~ns, 170--173 at 87.5~ns.
\item Pedestal noise at 43~V, 25~ns: 0.068--0.075 photoelectrons, 0.033--0.038 coherent,
  0.059--0.065 incoherent.
\item Threshold: one photoelectron is 22 DAC units at 43~V, 12.5 at 41~V; edge of the noise at
  190 and 182.
\end{enumerate}
\end{frame}
